\section{A certification frame for one-dimensional threshold management}\label{sec:frame}

\subsection{Objects}\label{sec:objects}

The state is a scalar $x\in X=[x_{\mathrm{lo}},x_{\mathrm{hi}}]$ (spawning biomass in kt; index-well head in ft). A \emph{threshold} $\LRP\in X$ defines the safe set $\safe=[\LRP,x_{\mathrm{hi}}]$; the ceiling $x_{\mathrm{hi}}$ is the top of the model's declared domain, and the state is clamped there rather than failed. A \emph{rule} (policy) is a map $\pi:X\to U$ from state to control, piecewise constant with finitely many jump points in all cases treated here. The one-step dynamics are
\begin{equation}\label{eq:dyn}
x_{t+1}=\Phi(x_t,\pi(x_t),w_t),
\end{equation}
where $w_t$ is a scalar disturbance---the residual of the fitted map for cod, the annual recharge for the aquifer---and $\Phi$ is non-decreasing in $w$. An \emph{uncertainty class} is a set of admissible disturbance sequences; the classes used here are \emph{floor classes} $\{w_t\ge \underline w\ \forall t\}$ with $\underline w$ the minimum, 5th or 10th percentile of the fit-window disturbances (``worst'', ``q05'', ``q10''). The \emph{floor map} is $F_\pi(x)=\Phi(x,\pi(x),\underline w)$.

\begin{definition}[closed-loop kernels]\label{def:kernel}
For a rule $\pi$ and a class $\mathcal W$, the $T$-step closed-loop robust kernel is
$W_T(\pi,\mathcal W)=\{x\in\safe:\ x_t\in\safe\ \text{for } t=0,\dots,T \text{ for every } (w_t)\in\mathcal W\}$, and $W_\infty=\bigcap_T W_T$. Its lower boundary is $b_T=\inf W_T$ ($b_T=\infty$ if $W_T=\emptyset$).
\end{definition}

\begin{lemma}[floor reduction]\label{lem:floor}
If $\Phi$ is non-decreasing in $w$ and $\mathcal W$ is a floor class, then $W_T(\pi,\mathcal W)$ equals the deterministic $T$-step viability kernel of $\safe$ under the floor map $F_\pi$, i.e.\ the sets generated by
\begin{equation}\label{eq:rec}
W_0=\safe,\qquad W_{t+1}=\{x\in\safe:\ F_\pi(x)\in W_t\}.
\end{equation}
\end{lemma}
\begin{proof}
Fix $x$ and a disturbance sequence in the class. By monotonicity in $w$ and induction, the realised state dominates the floor orbit, $x_t\ge F_\pi^{\,t}(x)$, provided the floor orbit stays in the region on which $F_\pi$ is monotone in $x$ (verified below for each map). Hence the floor sequence is the worst admissible one, and $x$ is $T$-step safe for every admissible sequence iff it is $T$-step safe under the floor map, which is \eqref{eq:rec}.
\end{proof}

The recursion \eqref{eq:rec} is exact and cheap in one dimension. If $F_\pi$ is strictly increasing on each piece of $\pi$, the preimage of an interval under a piece is an interval obtained from the branch inverse, so every $W_t$ is a finite union of closed intervals and \eqref{eq:rec} is a finite computation (the registered runners and the independent core of Section~\ref{sec:verification} both implement it this way).

\subsection{The kernel of a monotone map above a threshold}

\begin{theorem}[kernel structure]\label{thm:kernel}
Let $F:\safe\to\R$ be continuous and strictly increasing with $F(x_{\mathrm{hi}})\le x_{\mathrm{hi}}$, and let $W_t$ be generated by \eqref{eq:rec}. Then:
\begin{enumerate}[label=(\alph*)]
\item each $W_t$ is a closed interval $[b_t,x_{\mathrm{hi}}]$ or empty, with $b_0=\LRP$ and $b_{t+1}=\max\{\LRP,F^{-1}(b_t)\}$, where $F^{-1}(y)=\inf\{x\in\safe:F(x)\ge y\}$ and $F^{-1}(y)=\infty$ if $F(x_{\mathrm{hi}})<y$; the sequence $b_t$ is non-decreasing;
\item if $F(\LRP)\ge\LRP$ then $W_t=\safe$ for all $t$ (\emph{whole safe set});
\item if $F(\LRP)<\LRP$ and $F$ has a fixed point in $(\LRP,x_{\mathrm{hi}}]$, let $x_f$ be the smallest; then $b_t\uparrow x_f$ with $b_t<x_f$ for every finite $t$, and $W_\infty=[x_f,x_{\mathrm{hi}}]$ (\emph{half-line});
\item if $F(\LRP)<\LRP$ and $F$ has no fixed point in $(\LRP,x_{\mathrm{hi}}]$, then $W_T=\emptyset$ for all $T\ge\Tempty=\min\{T:F^{T}(x_{\mathrm{hi}})<\LRP\}<\infty$ (\emph{empty beyond a finite horizon}).
\end{enumerate}
\end{theorem}
\begin{proof}
(a) By induction. If $W_t=[b_t,x_{\mathrm{hi}}]$ then $F(x)\in W_t$ iff $b_t\le F(x)\le x_{\mathrm{hi}}$; the upper inequality holds on $\safe$ because $F(x)\le F(x_{\mathrm{hi}})\le x_{\mathrm{hi}}$, and by strict monotonicity the lower one holds iff $x\ge F^{-1}(b_t)$. Monotonicity of $b_t$: $b_1\ge\LRP=b_0$, and $b_t\ge b_{t-1}$ implies $F^{-1}(b_t)\ge F^{-1}(b_{t-1})$. (b) $F(\LRP)\ge\LRP$ gives $F^{-1}(\LRP)\le\LRP$, so $b_1=\LRP$ and inductively $b_t=\LRP$. (c) On $[\LRP,x_f)$ we have $F(x)<x$ (continuity, $F(\LRP)<\LRP$, no fixed point), so an orbit started there decreases strictly and, having no fixed point to converge to in $[\LRP,x_f)$, leaves $\safe$ in finite time; hence $b_\infty\ge x_f$. Conversely $[x_f,x_{\mathrm{hi}}]$ is forward invariant ($F(x_f)=x_f$, $F$ increasing, $F(x_{\mathrm{hi}})\le x_{\mathrm{hi}}$), so $b_t\le x_f$ for all $t$. Strictness: $b_t<x_f$ implies $b_{t+1}=\max\{\LRP,F^{-1}(b_t)\}<F^{-1}(x_f)=x_f$. (d) Since $F$ is increasing, $F^{t}(x)\le F^{t}(x_{\mathrm{hi}})$ for all $x\in\safe$, so $W_T\neq\emptyset$ iff $x_{\mathrm{hi}}\in W_T$ iff $F^{t}(x_{\mathrm{hi}})\ge\LRP$ for $t\le T$. The orbit of $x_{\mathrm{hi}}$ is non-increasing and cannot converge inside $\safe$ (no fixed point there), so it crosses $\LRP$ at a finite $\Tempty$.
\end{proof}

The theorem says that a monotone one-dimensional threshold problem has exactly three outcomes, and that the outcome is decided by two scalar facts: the sign of $F(\LRP)-\LRP$ and the existence of a fixed point above $\LRP$. The two corollaries record what those facts are for the two model families used here.

\begin{corollary}[surplus-production stock]\label{cor:cod}
Let $F_C(S)=S+g(S)-C-|e|$ with $g$ continuous, $g$ increasing on $[0,S_{\mathrm{pk}}]$ and decreasing beyond, $\gmax=g(S_{\mathrm{pk}})$, $\LRP<S_{\mathrm{pk}}$, $F_C$ strictly increasing on $\safe$ and $g(x_{\mathrm{hi}})\le C+|e|$. Define
\[
\Cstar=g(\LRP)-|e|,\qquad \Cvac=\gmax-|e|.
\]
Then for a constant catch $C\ge0$: $W_\infty=\safe$ iff $C\le\Cstar$; $\safe\supsetneq W_\infty=[s_-(C+|e|),x_{\mathrm{hi}}]\neq\emptyset$ iff $\Cstar<C\le\Cvac$, where $s_-(\ell)$ is the smaller root of $g(S)=\ell$; and $W_T=\emptyset$ beyond a finite horizon iff $C>\Cvac$. The $T$-step boundaries follow from (a) with the branch inverse of $F_C$. For the Schaefer form $g(S)=rS(1-S/K)$: $\gmax=rK/4$, $s_\pm(\ell)=\tfrac K2\bigl(1\pm\sqrt{1-4\ell/(rK)}\bigr)$, $F_C'(S)=1+r(1-2S/K)$, and $F_C'(\LRP)>1$ iff $K>2\LRP$. For a proportional rule $C=\phi\,g(S)$ on $\safe$ the same statements hold with $g$ replaced by $(1-\phi)g$: whole safe set iff $\phi\le 1-|e|/g(\LRP)$, non-empty iff $\phi\le1-|e|/\gmax$.
\end{corollary}

\begin{corollary}[affine aquifer map]\label{cor:edw}
Let $F(H)=aH+b$ with $0<a<1$, $b=\alpha+\beta\underline R+\gamma P$, and attractor $\Hstar=b/(1-a)\le x_{\mathrm{hi}}$. Then $W_\infty=\safe$ iff $\Hstar\ge\LRP$; otherwise there is no fixed point above $\LRP$,
\[
b_T=\Hstar+\frac{\LRP-\Hstar}{a^{T}},\qquad \Tempty=\Bigl\lfloor\frac{\ln\bigl((x_{\mathrm{hi}}-\Hstar)/(\LRP-\Hstar)\bigr)}{\ln(1/a)}\Bigr\rfloor+1 .
\]
The half-line regime does not occur for a contraction: the kernel of a threshold above the attractor is the whole safe set or empty in finite time.
\end{corollary}
\begin{proof}
Both corollaries instantiate Theorem~\ref{thm:kernel}: $F_C(\LRP)\ge\LRP$ iff $g(\LRP)\ge C+|e|$; fixed points of $F_C$ solve $g(S)=C+|e|$, which has a root above $\LRP$ iff $C+|e|\le\gmax$ (then $s_-(C+|e|)>\LRP$ because $g(\LRP)<C+|e|$ and $g$ increases up to $S_{\mathrm{pk}}$). For the affine map the unique fixed point is $\Hstar$; $F(\LRP)\ge\LRP$ iff $\Hstar\ge\LRP$; $b_T$ follows from $F^{-1}(y)=(y-b)/a$ by induction, and $\Tempty$ is the first $T$ with $b_T>x_{\mathrm{hi}}$.
\end{proof}

\begin{proposition}[stage-at-threshold principle]\label{prop:stage}
Let $\pi$ be piecewise constant with jump points $s_1<s_2<\dots$ and let $F_\pi$ be strictly increasing on each piece. (i) The kernels $W_T(\pi)$ of a threshold $\LRP$ depend on $\pi$ only through its values on $\safe=[\LRP,x_{\mathrm{hi}}]$; two rules that agree on $\safe$ have identical kernels of $\LRP$ at every horizon. (ii) $W_\infty(\pi)=\safe$ iff $F_\pi(\LRP)\ge\LRP$ and $F_\pi(s_j^{+})\ge\LRP$ at every jump point $s_j\in\safe$, i.e.\ iff the stage active at the threshold holds it and no higher stage dips below it. (iii) In particular, if $\pi$ coincides on $\safe$ with a constant control $u_\LRP$, then $W_T(\pi)=W_T(u_\LRP)$ for all $T$, and the stages prescribed below $\LRP$ are irrelevant to the certification of $\LRP$.
\end{proposition}
\begin{proof}
(i) is immediate from \eqref{eq:rec}, which only evaluates $F_\pi$ on $\safe$. (ii) $\safe$ is forward invariant iff $\inf_{x\in\safe}F_\pi(x)\ge\LRP$; on each piece the infimum is attained at the left end by monotonicity. (iii) is (i).
\end{proof}

The proposition is elementary, but it is the statement that most often contradicts the design intent of staged rules. A cascade that cuts harder as the state falls further below a trigger protects the \emph{lower} lines it draws; it does nothing for the trigger itself that the first stage does not already do. In Section~\ref{sec:cod} this makes the Stage-I rule, the declared cascade and a flat 60-kt cap indistinguishable at the LRP; in Section~\ref{sec:edwards} it makes the Stage-I rule tie the flat 80\% cap at 618 ft and Critical Period Management tie the flat 60\% cap.

\subsection{Model defect: exact robust kernel and the Lipschitz certificate}\label{sec:defect}

Floor classes bound the disturbance; they do not bound the error of the fitted map itself. The registered protocols declared a uniform defect $\varepsilon$ equal to the largest absolute fit-window residual and certified against the class $\mathcal D_\varepsilon$ of all maps $F+\delta_t$ with $\sup_x|\delta_t(x)|\le\varepsilon$, time-varying and state-dependent, applied on top of the floor class.

\begin{theorem}[exact bounded-defect kernel]\label{thm:exact}
For a floor map $F$ (clamped at $x_{\mathrm{hi}}$) and the defect class $\mathcal D_\varepsilon$, the $T$-step robust kernel of $\safe$ is generated by
\begin{equation}\label{eq:exact}
W^{\varepsilon}_0=\safe,\qquad W^{\varepsilon}_{t+1}=\bigl\{x\in\safe:\ [\,F(x)-\varepsilon,\ \min\{F(x)+\varepsilon,x_{\mathrm{hi}}\}\,]\subset W^{\varepsilon}_t\bigr\}.
\end{equation}
If $F$ is strictly increasing on each piece, every $W^{\varepsilon}_t$ is a finite union of intervals computable from the branch inverses; for a single increasing branch with the ceiling non-binding, $W^{\varepsilon}_T=[b^{\varepsilon}_T,x_{\mathrm{hi}}]$ with $b^{\varepsilon}_0=\LRP$ and $b^{\varepsilon}_{t+1}=\max\{\LRP,F^{-1}(b^{\varepsilon}_t+\varepsilon)\}$.
\end{theorem}
\begin{proof}
A state is $(t{+}1)$-step robustly safe iff it is safe now and every successor the class allows is $t$-step robustly safe; the set of allowed successors of $x$ is the (clamped) interval $[F(x)-\varepsilon,F(x)+\varepsilon]$ because the defect is an arbitrary bounded perturbation at every step \cite{BertsekasRhodes1971,Aubin1991}. Nothing is lost in \eqref{eq:exact}: it is the dynamic-programming characterisation of the kernel, not a bound on it. The interval form follows as in Theorem~\ref{thm:kernel}(a).
\end{proof}

\begin{proposition}[Lipschitz certificate]\label{prop:lip}
Let $a\ge\sup_{\safe}F'$ and $r_T=\varepsilon\sum_{k=0}^{T-1}a^{k}$, i.e.\ $r_T=\varepsilon(a^{T}-1)/(a-1)$ for $a\neq1$. Write $W_T(\LRP')$ for the nominal $T$-step kernel of a threshold $\LRP'$ under \eqref{eq:rec}.
\begin{enumerate}[label=(\roman*)]
\item (inner approximation) $W_T(\LRP+r_T)\subset W^{\varepsilon}_T$ for every $T$.
\item (affine case) If $F(x)=ax+b$ then $b^{\varepsilon}_T=\max\{\LRP,F^{-T}(\LRP+r_T)\}$: the exact boundary is the Lipschitz boundary without the inflation of the initial state. The two certificates coincide whenever the backward recursion binds.
\item (expansive concave case) If $F'$ is strictly decreasing on $\safe$ (as for the Schaefer map) and $W^{\varepsilon}_T\neq\emptyset$ with $T\ge2$, the inclusion in (i) is strict.
\end{enumerate}
\end{proposition}
\begin{proof}
(i) Let $x\in W_T(\LRP+r_T)$ and let $x_t$ be a realised orbit under some $F+\delta_t$. With $d_t=|x_t-F^{t}(x)|$ one has $d_{t+1}\le a\,d_t+\varepsilon$, so $d_t\le r_t\le r_T$ and $x_t\ge F^{t}(x)-r_T\ge\LRP$ for $t\le T$. (ii) For an affine map $F^{-t}(y+a^{t}\varepsilon)=F^{-t}(y)+\varepsilon$, so by induction $F^{-1}\bigl(F^{-t}(\LRP+r_t)+\varepsilon\bigr)=F^{-(t+1)}(\LRP+r_t+a^{t}\varepsilon)=F^{-(t+1)}(\LRP+r_{t+1})$, which is the recursion of Theorem~\ref{thm:exact}. (iii) The exact recursion propagates the defect through $F^{-1}$ evaluated at $b^{\varepsilon}_t+\varepsilon>\LRP$, where the local slope is strictly below $a=F'(\LRP)$; the Lipschitz recursion charges the slope $a$ at every step, so for $T\ge2$ its boundary is strictly higher wherever both are finite.
\end{proof}

\begin{proposition}[finite certified horizons]\label{prop:finite}
Let $\Tcert=\max\{T:W^{\varepsilon}_T\neq\emptyset\}$ for a constant control.
\begin{enumerate}[label=(\roman*)]
\item (surplus-production stock) $\Tcert<\infty$ iff $C+|e|+\varepsilon>\gmax$, i.e.\ iff $C>\Cvac-\varepsilon$. If $\varepsilon\ge\Cvac$ no catch level, including zero, is certifiable indefinitely. Under the Lipschitz certificate, in addition, $\Tcert\le\bigl\lfloor\ln\bigl(1+(x_{\mathrm{hi}}-\LRP)(a-1)/\varepsilon\bigr)/\ln a\bigr\rfloor$ for every rule.
\item (affine contraction) $\Tcert<\infty$ iff $\Hstar<\LRP+\varepsilon/(1-a)$; a rule is certifiable indefinitely iff its worst-floor attractor clears the threshold by at least $r_\infty=\varepsilon/(1-a)$.
\end{enumerate}
\end{proposition}
\begin{proof}
(i) By Theorem~\ref{thm:exact} the exact boundary obeys $b_{t+1}=F^{-1}(b_t+\varepsilon)$, which has a fixed point iff $F(x)-\varepsilon=x$ has a solution above $\LRP$, i.e.\ $g(x)=C+|e|+\varepsilon$, possible iff the right side is at most $\gmax$. Without a fixed point $b_t$ increases without bound and passes $x_{\mathrm{hi}}$ at a finite $T$. The Lipschitz bound follows from $W_T(\LRP+r_T)=\emptyset$ once $\LRP+r_T>x_{\mathrm{hi}}$. (ii) For the affine map, $b^{\varepsilon}_T=\Hstar+(\LRP+r_T-\Hstar)/a^{T}$ when the recursion binds; since $r_T\uparrow r_\infty$ and $a^{-T}\to\infty$, $b^{\varepsilon}_T$ exceeds $x_{\mathrm{hi}}$ for large $T$ iff $\LRP+r_\infty-\Hstar>0$; if $\Hstar\ge\LRP+r_\infty$ the recursion never binds and $W^{\varepsilon}_T\supset[\LRP+r_T,x_{\mathrm{hi}}]$ for all $T$.
\end{proof}

\subsection{Securability, retention and the frozen protocol}

\begin{theorem}[securability margin]\label{thm:margin}
Suppose the control enters the floor map monotonically and let $F_0$ be the floor map under the most protective admissible control (zero catch; zero pumping). A threshold $\LRP$ admits a rule with $W_\infty=\safe$ iff $F_0(\LRP)\ge\LRP$. For the surplus-production stock this reads $g(\LRP)\ge|e|$ ($\Cstar\ge0$); for the affine aquifer map it reads $\Hstar_0=(\alpha+\beta\underline R)/(1-a)\ge\LRP$, and the minimal flat cut securing $\LRP$ is $1-\rho^{*}$ with
$\rho^{*}=\bigl(\LRP(1-a)-\alpha-\beta\underline R\bigr)/(\gamma\Pbar)$.
\end{theorem}
\begin{proof}
If $F_0(\LRP)<\LRP$ then $F_\pi(\LRP)\le F_0(\LRP)<\LRP$ for every admissible rule, and Theorem~\ref{thm:kernel}(c,d) excludes the whole safe set. If $F_0(\LRP)\ge\LRP$ the most protective rule attains it by Theorem~\ref{thm:kernel}(b). The cut formula solves $\Hstar(\rho\Pbar)=\LRP$.
\end{proof}

A \emph{reading} is a pair (floor class, threshold). Following the registered protocols, a staged rule is \emph{retained} against the flat family at a reading if (a) it is at least as protective as business-as-usual at every horizon there, and (b) it supplies at least as much control (catch; pumping) as every flat cap that is at least as protective as itself. Clause (a) is a dominance condition on kernels, clause (b) a supply comparison on the observed state history; both are decided, not estimated.

The protocols fix, before any kernel is computed: the data extracts and their provenance; the fit window and the out-of-sample audit window; the model form and the estimator; the residual convention; the floor classes and the quantile convention (linear interpolation); the defect $\varepsilon$ as the largest absolute fit-window residual with the out-of-sample maximum reported as an audit; the thresholds; the ceiling; the policy family; the horizon grid; and the retention rule. Everything downstream is a theorem about those declarations.
